\documentclass{article}
\usepackage{amsmath,amsthm,amsfonts,amssymb}

\newtheorem{claim}{Claim}

\title{Fields and Galois Theory}

\author{Aditya Padiyar}

\date{August 2026}

\begin{document}

\maketitle

\section{Fields}

\begin{claim}
  There is a degree 4 extension of $\mathbb{Q}$ with exactly one intermediate subfield
\end{claim}

\begin{proof}
  Consider the fields $\mathbb{Q} \subset \mathbb{Q}(\alpha^2) \subset \mathbb{Q}(\alpha)$, where $\alpha = 2^\frac{1}{4}$ 
\end{proof}

\begin{claim}
  There is a degree 4 extension with infinitely many intermediate subfields
\end{claim}

\begin{proof}
  Consider the fields $F \subset F(\alpha +c \beta) \subset \frac{R}{I}$ where $F = \mathbb{F}_2(x,y)$, $R=F[s,t], I = (s^2-x,t^2-y), \alpha = s + I, \beta = t + I, c \in F$ 
\end{proof}

\begin{claim}
  Let $\alpha = \sqrt{2}$ and $\beta > 0$ be such that $\beta^3 - 5 = 0$. Then $[\mathbb{Q}(\alpha + \beta): \mathbb{Q}]=6$
\end{claim}

\begin{proof}
  Use the fact that a $\mathbb{Q}$-basis for $\mathbb{Q}(\alpha,\beta)$ is given by \[\{1,\beta,\beta^2,\alpha,\alpha\beta,\alpha\beta^2\}\]
\end{proof}

\begin{claim}
  Suppose we have $F \subset R \subset K$, where $F,K$ are fields, $R$ is a ring, and $K$ is algebraic over $F$. Then $R$ is a field
\end{claim}

\begin{proof}
  Let $r \neq 0 \in R$. The map $M$ from $F[r] \to F[r]$ given by $x \to rx$ is $F$-linear and injective. Since $F[r]$ is finite dimensional over $F$, $M$ must be surjective as well.
\end{proof}

\begin{claim}
  Suppose $F \subset K$ is a finite extension and $\alpha_1 \in K$. Consider the $F$-linear map $T$ from $K \to K$ given by $x \to \alpha_1 x$. 
  Then $T$ has a block diagonal matrix with respect to some $F$-basis of $K$
\end{claim}

\begin{proof}
  We let $K=F(\alpha_1,\dots,\alpha_n)$, so that a basis of $K$ is given by $\alpha_1^{j_1}\alpha_2^{j_2}\dots \alpha_n^{j_n}$. 
  Let $p(x) \in F[x]$ be the irreducible polynomial of $\alpha$ and $C$ its companion matrix. 
  The matrix of $M$ is then $\text{diag}(C,C,\dots,C)$
\end{proof}

\begin{claim}
  Let $\alpha = \text{cos}(\frac{2\pi}{7})$. Then $[\mathbb{Q}(\alpha):\mathbb{Q}]=3$
\end{claim}

\begin{proof}
  Let $w = e^{i\frac{2\pi}{7}}$. Then we have $w^2-2\alpha w +1 =0$
\end{proof}

\begin{claim}
  Suppose $f(x),g(x) \text{ are coprime in }F[x]$ with $\text{deg } f < \text{deg } g=n$. Let $\alpha=\frac{f(x)}{g(x)} \in F(x)$. Then $[F(x):F(\alpha)]=n$
\end{claim}

\begin{proof}
  Use Gauss's lemma
\end{proof}

\begin{claim}
  Suppose that every polynomial in $F[x]$ has a root in a separable algebraic extension $K$ of $F$. Then $K$ is algebraically closed.
\end{claim}

\begin{proof}
  Suppose $f(x)=(x-\alpha_1)\dots(x-\alpha_n) \in F[x]$ where $\alpha_i \in \overline{K}$. 
  By the primitive element theorem we have $F(\alpha_1, \dots,  \alpha_n)= F(\beta)$. 
  Let $p(x) \in F[x]$ be the  minimal polynomial of $\beta$ over $F$. 
  By the hypothesis, $p(\alpha)=0$ with $\alpha \in K$. 
  Hence we have an isomorphism $\phi:F(\beta) \to F(\alpha)$ that fixes $F$ with $\phi(\beta)=\alpha$. 
  Hence $\alpha_i \in K$
\end{proof}

\section{Some Galois groups}

\begin{claim}
  Let $K=\mathbb{Q}(\sqrt{2},\sqrt{3},\sqrt{5})$. Then the following hold:
  \begin{itemize}
    \item $K$ is a degree $8$ Galois extension over $\mathbb{Q}$
    \item $\text{Aut}(K/\mathbb{Q})$ is isomorphic to $(\frac{\mathbb{Z}}{2\mathbb{Z}})^3$
    \item there are 16 intermediate fields
  \end{itemize}
\end{claim}

\begin{proof}
  Consider the following:
  \begin{itemize}
    \item This is easy to verify
    \item Since every element of $\text{Aut}(K/\mathbb{Q})$ has degree at most 2 the only possible group is $(\frac{\mathbb{Z}}{2\mathbb{Z}})^3$
    \item this is because $(\frac{\mathbb{Z}}{2\mathbb{Z}})^3$ has 16 subgroups
  \end{itemize}
\end{proof}


\begin{claim}
  Let $K=\mathbb{Q}(\alpha,i)$, where $\alpha=2^\frac{1}{4}$. Then the following hold:
  \begin{itemize}
    \item $K$ is a degree 8 Galois extension over $\mathbb{Q}$
    \item $\text{Aut}(K/\mathbb{Q})$ is isomorphic to the dihedral group with 8 elements
  \end{itemize}
\end{claim}

\begin{proof}
  Consider the following:
  \begin{itemize}
    \item this is easy to verify
    \item Define the following elements of $\text{Aut}(K/\mathbb{Q})$:
          \begin{itemize}
            \item let $\sigma$ be such that $\sigma(\alpha)=\alpha i$ and $\sigma(i)=i$
            \item let $\tau$ be such that $\tau(\alpha)=\alpha$ and $\tau(i)=-i$ 
          \end{itemize}
          Then we have:
          \begin{itemize}
            \item $\sigma^4=1$
            \item $\tau^2=1$ 
            \item $\sigma \tau = \tau \sigma^3$
          \end{itemize}
  \end{itemize}
\end{proof}

\begin{claim}
  Let $K=\mathbb{Q}(\omega_7)$, where $\omega_7 = e^{i\frac{2\pi}{7}}$. Then we have:
  \begin{itemize}
    \item $K$ is a degree 6 Galois extension of $\mathbb{Q}$
    \item $\text{Aut}(K/\mathbb{Q})$ is isomorphic to $\frac{\mathbb{Z}}{6\mathbb{Z}}$
  \end{itemize}
\end{claim}

\begin{proof}
  Consider the following:
  \begin{itemize}
    \item The cyclotomic polynomial is irreducible
    \item the element of $\text{Aut}(K  / \mathbb{Q})$ that sends $\omega_7 \to \omega_7^3$ has order 6
  \end{itemize}
\end{proof}

\begin{claim}
  Let $K=\mathbb{Q}(e^{i\frac{\pi}{4}})$. Then the following hold:
  \begin{itemize}
    \item $K$ is a degree 4 Galois extension over $\mathbb{Q}$
    \item $\text{Aut}(K/\mathbb{Q})$ is isomorphic to the Klein four group
  \end{itemize}
\end{claim}

\begin{proof}
  Consider the following:
  \begin{itemize}
    \item this is the splitting field of $x^4+1$
    \item there are no order 4 elements
  \end{itemize}
\end{proof}

\begin{claim}
  Let $p(x)=x^5-4x+2 \in \mathbb{Q}[x]$ and $K$ be its splitting field. Then the following hold:
  \begin{itemize}
    \item $p$ is irreducible over $\mathbb{Q}$
    \item $p$ has exactly 3 real roots
    \item $K$ has degree 120 over $\mathbb{Q}$ 
    \item $G = \text{Aut}(K/\mathbb{Q})$ is isomorphic to $S_5$
  \end{itemize}
\end{claim}

\begin{proof}
  Consider the following:
  \begin{itemize}
    \item This follows from the Eisenstein criterion
    \item this is easy to check
    \item Consider the following:
          \begin{itemize}
            \item $G$ acts faithfully on the roots of $p$. Since $p$ is irreducible the action is transitive.
            \item We see that there is an intermediate subfield $L$ of degree 5 over $\mathbb{Q}$. Hence 5 must divide $|G|$.
            \item There must exist a 5 cycle in $G$
            \item complex conjugation is a two cycle in $G$
            \item any 5 cycle and any two cycle generate $S_5$
          \end{itemize} 
  \end{itemize}
\end{proof}

\begin{claim}
  Let $p(x)=x^3-3x+1 \in \mathbb{Q}[x]$ and $K$ be its splitting field. Then the following hold:
  \begin{itemize}
    \item $p$ is irreducible over $\mathbb{Q}$
    \item $p$ has exactly 3 real roots
    \item $\text{Aut}(K/\mathbb{Q})$ is isomorphic to $\frac{\mathbb{Z}}{3\mathbb{Z}}$
  \end{itemize}
\end{claim}

\begin{proof}
  Consider the following:
  \begin{itemize}
    \item reduce mod 2
    \item easy to verify
    \item Consider the following:
          \begin{itemize}
            \item $G$ acts faithfully and transitively on the roots of $p$
            \item the discriminant $\delta^2 = 81$, which means $\delta \in \mathbb{Q}$
            \item this forces every permutation to be even
          \end{itemize}
  \end{itemize}
\end{proof}

\begin{claim}
  Let $K=\mathbb{Q}(\alpha,\omega_5)$, where $\alpha=2^\frac{1}{5}$. Then $K$ is a degree 20 Galois extension of $\mathbb{Q}$
\end{claim}

\begin{proof}
  Let $p(x)=x^5-2$. Consider the following:
  \begin{itemize}
    \item $K$ is the splitting field of $p$ and hence a Galois extension
    \item $p$ is irreducible by Eisenstein 
    \item $w_5$ has degree 4 over $\mathbb{Q}$
    \item hence both 4 and 5 must divide $[K:\mathbb{Q}]$, which means $[K:\mathbb{Q}] \geq 20$
    \item the diamond rule says that $[K:\mathbb{Q}] \leq 20$
  \end{itemize}
\end{proof}

\begin{claim}
  Let $K=\mathbb{Q}(\alpha,\omega_6)$, where $\alpha=2^\frac{1}{6}$. Then the following hold:
  \begin{itemize}
    \item $K$ is a degree 12 Galois extension of $\mathbb{Q}$
    \item $\text{Aut}(K/\mathbb{Q})$ is isomorphic to $\frac{\mathbb{Z}}{2\mathbb{Z}} \times S_3$
  \end{itemize}
\end{claim}

\begin{proof}
  Consider the following:
  \begin{itemize}
    \item $K$ is the splitting field of the irreducible polynomial $x^6-2$ which makes it a Galois extension.
    \item check that $K$ has degree 12 over $\mathbb{Q}$
    \item define the following elements of $\text{Aut}(K/\mathbb{Q})$
    \begin{itemize}
      \item $\sigma_1(\alpha) = \alpha w^2$ and $\sigma_1(\omega)=\omega$
      \item $\sigma_2(\alpha) = \alpha w^3$ and $\sigma_2(\omega)=\omega$
      \item $\text{conj}$ which is complex conjugation
    \end{itemize}
    \item let $H$ be the subgroup generated by ($\sigma_1$, $\text{conj}$)
    \item check that $H$ is isomorphic to $S_3$
    \item let $K$ be the subgroup generated by $\sigma_2$ which is isomorphic to $\frac{\mathbb{Z}}{2\mathbb{Z}}$
    \item we can now use Proposition 2.11.4 of Artin 
  \end{itemize}
\end{proof}

\begin{claim}
  There is a field extension $K$ of $\mathbb{Q}$ such that $\text{Aut}(K / \mathbb{Q})$ is isomorphic to $\frac{\mathbb{Z}}{n\mathbb{Z}}$
\end{claim}

\begin{proof}
  Consider the following:
  \begin{itemize}
    \item choose $p$ such that $p$ has residue 1 mod $n$
    \item let $L=\mathbb{Q}(\omega_p)$ which is a Galois extension of $\mathbb{Q}$ of degree $p-1$
    \item let $G$ be the Galois group 
    \item check that $G$ is isomorphic to $\mathbb{F}_p^\times$ which is isomorphic to $\frac {\mathbb{Z}}   {   (p-1) \mathbb{Z} }  $
    \item let $H$ be a subgroup of $G$ of index $n$
    \item set $K=L^H$. Then $\text{Aut}(K / \mathbb{Q}) \cong \frac{G}{H} \cong \frac{\mathbb{Z}}{n\mathbb{Z}}$
  \end{itemize}
\end{proof}

\section{Abstract Galois extensions}

Suppose that $K$ is a finite Galois extension of $F$ with Galois group $G$

\begin{claim}
  Suppose that:
  \begin{itemize}
    \item $H_1$ and $H_2$ be subgroups of $G$
    \item $K_1=K^{H_1}$ and $K_2=K^{H_2}$
  \end{itemize}

  Then the following hold:
  \begin{itemize}
    \item $K_1 \cap K_2 = K^{H_3}$ where $H_3$ is the subgroup generated by $H_1$ and $H_2$
    \item $K^{H_1 \cap H_2} = L$ where $L$ is the subfield generated by $K_1$ and $K_2$
  \end{itemize}
\end{claim}

\begin{proof}
  this is easy
\end{proof}

\begin{claim}
  If $H$ is a subgroup of $G$ then there is $\alpha \in K$ whose stabilizer under the action of $G$ is $H$
\end{claim}

\begin{proof}
  Let $L=K^H$. By the primitive element theorem $L=F(\alpha)$. Then the stabilizer is $\text{Aut}(K/L)$ which is $H$.
\end{proof}

\begin{claim}
  Suppose $G$ is isomorphic to $\frac{\mathbb{Z}}{n\mathbb{Z}}$ and $d$ divides $n$. Then there is a unique intermediate field that has degree $d$ over $F$
\end{claim}

\begin{proof}
  there is a unique subgroup of index $d$
\end{proof}

\end{document}
